\documentclass[11pt]{article}
\usepackage[margin=1in]{geometry}
\usepackage[T1]{fontenc}
\usepackage{lmodern}
\usepackage{amsmath,amssymb,booktabs,longtable,array,enumitem,hyperref,xcolor}
\hypersetup{colorlinks=true,linkcolor=blue,urlcolor=blue}
\setlength{\emergencystretch}{3em}
\setlength{\tabcolsep}{3.5pt}
\newcolumntype{L}[1]{>{\raggedright\arraybackslash}p{#1}}
\newcommand{\E}{\mathbb{E}}
\newcommand{\Cov}{\operatorname{Cov}}
\newcommand{\Var}{\operatorname{Var}}
\newcommand{\code}[1]{\nolinkurl{#1}}

% Replace these fields before an audit begins.
\newcommand{\AuditPackageID}{PACKAGE-ID}
\newcommand{\AuditPackageName}{Package scientific-contract audit}
\newcommand{\AuditTier}{A}
\newcommand{\AuditGoverningSHA}{TO-COMPLETE}
\newcommand{\AuditBranch}{TO-COMPLETE}
\newcommand{\AuditWorktree}{TO-COMPLETE}
\newcommand{\AuditCorePath}{TO-COMPLETE}
\newcommand{\AuditCoreSHA}{TO-COMPLETE}
\newcommand{\AuditCoreFreezeCommit}{TO-COMPLETE}
\newcommand{\AuditDate}{TO-COMPLETE}

\title{\AuditPackageName\\
\large Citlali scientific-contract audit: \AuditPackageID}
\author{Independent package auditor}
\date{\AuditDate}

\begin{document}
\maketitle

\begin{abstract}
State the claimed estimator, the exact implementation assessed, the strongest
conclusion supported by the evidence, and the production restriction in no
more than one paragraph. This template is intentionally compact; use the
optional appendices only when the evidence requires them.
\end{abstract}

\section{Executive summary and bounded scope}

\begin{description}[style=nextline]
\item[Package] \AuditPackageID: \AuditPackageName.
\item[Tier] \AuditTier.
\item[Included transformation] State the single bounded estimator or
engineering contract.
\item[Explicit exclusions] Name adjacent packages and implementation paths
that are not being audited.
\item[Claimed result] State the audit conclusion without implying production
approval.
\item[Consumer policy] State the allowlist, restricted uses, and fail-closed
uses.
\end{description}

\section{Authority and independence record}

\begin{longtable}{L{0.27\linewidth}L{0.65\linewidth}}
\toprule
Field & Exact value\\
\midrule
Governing source SHA & \code{\AuditGoverningSHA}\\
Audit branch & \code{\AuditBranch}\\
Audit worktree & \code{\AuditWorktree}\\
Independent-core bytes & \code{\AuditCorePath}\\
Independent-core SHA-256 & \code{\AuditCoreSHA}\\
Freeze commit & \code{\AuditCoreFreezeCommit}\\
First package-source inspection & Record timestamp, command or source, and
confirmation that it followed the freeze.\\
Prior implementation exposure & List unavoidable prior knowledge or write
``none known''; explain how its influence was bounded.\\
\bottomrule
\end{longtable}

Identify the project-level scientific authorities used before the freeze and
the implementation files quarantined until afterward. The digest must cover
the exact independent mathematical core, not this evolving whole document.

\section{Input contract and variable classification}

State scientific identity, shape, ordering, units, coordinate frame, index
base, validity, missing/non-finite policy, and upstream assumptions for every
input.

\begin{longtable}{L{0.14\linewidth}L{0.18\linewidth}L{0.15\linewidth}L{0.17\linewidth}L{0.24\linewidth}}
\toprule
Quantity & Identity and shape & Units/frame & Classification & Assumption or
conditioning\\
\midrule
\endfirsthead
\toprule
Quantity & Identity and shape & Units/frame & Classification & Assumption or
conditioning\\
\midrule
\endhead
\(x\) & Input data & complete & random & complete\\
\(\theta\) & Policy or parameter & complete & fixed, fitted, selected, or
data-derived & complete\\
\bottomrule
\end{longtable}

Use only these classifications unless a package-specific term is defined:
random, fixed, fitted, selected, data-derived, and conditioned upon.

\section{Independently claimed estimator}
\label{sec:independent-core}

This section, together with any explicitly named companion core file, is the
content frozen before implementation inspection. Select the applicable form
and give a specific rationale for the others being not applicable.

\subsection{Affine form}

\begin{align}
  \boldsymbol y &= A\boldsymbol x+\boldsymbol c,
  \label{eq:affine}\\
  \E[\boldsymbol y] &= A\E[\boldsymbol x]+\boldsymbol c,
  \label{eq:affine-mean}\\
  C_y &= A C_x A^{\mathsf T}.
  \label{eq:affine-cov}
\end{align}

Define every coefficient and state whether \(A\) and \(\boldsymbol c\) are
fixed or data-derived. If a fitted or selected quantity is conditioned upon,
derive the conditional result and identify the omitted unconditional terms.

\subsection{Nonlinear or data-dependent form}

\begin{equation}
  \boldsymbol y=F(\boldsymbol x,\boldsymbol\theta).
  \label{eq:nonlinear}
\end{equation}

State the distributional target, fitted/selected state, linearization or
resampling method, and conditions under which uncertainty claims hold.

\subsection{Iterative or stateful form}

\begin{align}
  \boldsymbol s_{n+1} &= G(\boldsymbol s_n,\boldsymbol x),
  \label{eq:state-update}\\
  \boldsymbol y_n &= H(\boldsymbol s_n,\boldsymbol x).
  \label{eq:state-output}
\end{align}

Define initialization, absolute iteration identity, stopping rule, fixed point
or terminal estimator, restart state, and propagation of uncertainty between
iterations.

\section{Response and normalization}

Derive every applicable temporal, spatial, beam/template, extended-mode,
photometric, and iteration-to-iteration response. Distinguish DC, peak,
integrated, template, and matched-filter normalization. For a named template
\(\boldsymbol p\), a useful generic definition is
\begin{equation}
  R(\boldsymbol p)=\frac{\partial \E[\boldsymbol y]}
                         {\partial \alpha}
  \quad\hbox{for}\quad
  \E[\boldsymbol x]=\alpha\boldsymbol p+\boldsymbol\mu_0.
  \label{eq:template-response}
\end{equation}
State edge behavior, conditioning, and response uncertainty.

\section{Uncertainty and covariance}

Treat full covariance as the mathematical object and identify every diagonal,
independence, stationarity, or fixed-selection approximation. Address:

\begin{itemize}
\item formal variance and full covariance of the named estimator;
\item empirical or jackknife variance and its degrees of freedom;
\item calibration/systematic and response uncertainty;
\item covariance introduced by filtering, selection, fitting, or feedback;
\item the exact meaning and units of every inverse-variance product; and
\item what the stored uncertainty does not represent.
\end{itemize}

Never use inverse variance, support, validity, hits, coverage, confidence, or
response as synonyms.

\section{Support, validity, identity, and state}

Define support and invalidity operators separately from uncertainty. Record
units, frames, normalization, shapes, index bases, missing/non-finite behavior,
and detector/array/network/map identity.

\begin{longtable}{L{0.19\linewidth}L{0.19\linewidth}L{0.24\linewidth}L{0.26\linewidth}}
\toprule
State & Value represented & Resolution rule & Provenance/output\\
\midrule
Requested & Immutable accepted request & complete & complete\\
Effective & Context-free policy & complete & complete\\
Observation-resolved & Data-context choice & complete & complete\\
Realized & Applied/produced result & complete & complete\\
\bottomrule
\end{longtable}

\section{Implementation conformity trace}

Complete this section only after the independent-core freeze. Trace signal,
formal uncertainty, empirical/noise realizations, masks, simulations,
sequential/OpenMP paths, output writers, metadata, and consumers at the exact
governing SHA.

\begin{longtable}{L{0.27\linewidth}L{0.15\linewidth}L{0.23\linewidth}L{0.23\linewidth}}
\toprule
Source path & Lines/symbol & Contract equation & Conformity/evidence\\
\midrule
\endfirsthead
\toprule
Source path & Lines/symbol & Contract equation & Conformity/evidence\\
\midrule
\endhead
complete & complete & Equation~\eqref{eq:affine} or N/A & complete\\
\bottomrule
\end{longtable}

\section{Product and consumer matrix}

\begin{longtable}{L{0.13\linewidth}L{0.14\linewidth}L{0.15\linewidth}L{0.09\linewidth}L{0.18\linewidth}L{0.19\linewidth}}
\toprule
Input & Operator & Estimator or product & Units & Metadata and validity & Consumer and
authorization\\
\midrule
\endfirsthead
\toprule
Input & Operator & Estimator or product & Units & Metadata and validity & Consumer and
authorization\\
\midrule
\endhead
complete & complete & complete & complete & complete &
allowed, restricted, or fail-closed\\
\bottomrule
\end{longtable}

For each persisted product state extension/variable/table identity,
normalization, support, covariance assumptions, authoritative/derived status,
producer, current and recommended metadata, and applicable tests.

\section{Falsifiable validation matrix}

\begin{longtable}{L{0.22\linewidth}L{0.25\linewidth}L{0.20\linewidth}L{0.21\linewidth}}
\toprule
Evidence method & Claim tested & Exact fixture/artifact & Gate or N/A rationale\\
\midrule
\endfirsthead
\toprule
Evidence method & Claim tested & Exact fixture/artifact & Gate or N/A rationale\\
\midrule
\endhead
Analytic limiting cases & complete & complete & complete\\
Small deterministic fixtures & complete & complete & complete\\
Synthetic injections & complete & complete & complete\\
Blank controls & complete & complete & complete\\
Sequential/OpenMP equivalence & complete & complete & complete\\
Same-SHA Unity reduction & complete & complete & complete\\
Astronomical recovery & complete & complete & complete\\
\bottomrule
\end{longtable}

Acceptance tolerances are specified before looking at candidate results and
are derived from analytic precision, finite-sample behavior, repeatability,
or an explicit scientific decision. A required-data skip is not a pass.

\section{Findings, dependencies, and scientific decisions}

\begin{longtable}{L{0.15\linewidth}L{0.15\linewidth}L{0.08\linewidth}L{0.12\linewidth}L{0.38\linewidth}}
\toprule
Stable ID & Class & Priority & Basis/conf. & Finding, owner, and closure gate\\
\midrule
\endfirsthead
\toprule
Stable ID & Class & Priority & Basis/conf. & Finding, owner, and closure gate\\
\midrule
\endhead
\AuditPackageID-F001 & finding class & P0--P3 & derived, high & Describe finding,
owner, and closure gate.\\
\bottomrule
\end{longtable}

List every upstream package dependency by stable ID, the exact required fact,
its status, and the consequence while open. Keep implementation defects,
contract gaps, scientific-policy decisions, evidence gaps, and dependency gaps
separate.

\section{Verdict and production disposition}

\begin{description}[style=nextline]
\item[Contract status] TO-COMPLETE.
\item[Implementation status] TO-COMPLETE at \code{\AuditGoverningSHA}.
\item[Validation status] TO-COMPLETE.
\item[Production status] TO-COMPLETE, with an explicit scope.
\item[Verdict] Pending, retain, amend, split, supersede, or reject.
\item[Allowed consumers] TO-COMPLETE.
\item[Restricted/fail-closed consumers] TO-COMPLETE.
\item[Repair and re-audit trigger] TO-COMPLETE.
\end{description}

\section{Machine-readable ledger proposal}

Include a YAML fragment matching \code{doc/audits/audit-ledger.yaml}. It is a
proposal for coordinator review, not permission to modify the canonical ledger
from a parallel audit branch.

\appendix
\section{Optional extended source trace}

Use this appendix for longer code tables, derivations, or historical
comparisons. Delete it when it adds no review value.

\section{Optional returned evidence inventory}

List exact artifact paths/hashes, commands, configurations, comparator
versions/options, tolerances, skips, and provenance. Do not embed bulky
reduction products in Git.

\end{document}
